\documentclass[12pt,a4paper]{article}
\usepackage[utf8]{inputenc}
\usepackage{amsmath,amssymb,amsthm}
\usepackage{geometry}
\usepackage{setspace}
\usepackage{hyperref}
\usepackage{natbib}
\geometry{margin=1in}
\onehalfspacing

\newtheorem{definition}{Definition}
\newtheorem{lemma}{Lemma}
\newtheorem{proposition}{Proposition}
\newtheorem{remark}{Remark}
\newtheorem{corollary}{Corollary}
\newtheorem{assumption}{Assumption}

\title{Cherry-Picking under Screening:\\ Selective Disclosure in Pharmaceutical Approval}
\author{Xueyang}
\date{\today}

\begin{document}
\maketitle

% ============================================================
\section{Introduction}
% ============================================================

A pharmaceutical firm conducts a clinical trial to persuade the
FDA to approve a drug. The FDA does not know the firm's
type---that is, the set of experiments available to the
firm---and must design an approval rule to screen different
types. This screening problem is studied by Min (2023), where
the FDA observes both the experiment chosen and its outcome.

I modify Min's model in one respect: the experimental outcome
is the firm's private information, and the firm chooses whether
to disclose it. If the firm discloses, the FDA observes both
the experiment and the verified outcome. If the firm remains
silent, the FDA observes nothing---neither the experiment
conducted nor its result.

This information structure captures a key feature of
pharmaceutical regulation: when a firm publishes or submits
trial results, the data are verifiable (hard evidence); but
when a firm buries a trial, the regulator often does not know
the details of the unreported experiment. Empirical evidence
supports this asymmetry. A \emph{Science}/AAAS investigation
found that sponsors violated ClinicalTrials.gov reporting
requirements more than 55\% of the time, and the FDA levied
zero fines over seven years of mandatory reporting. A BMJ study
of high-risk cardiovascular devices found that approximately
half of the clinical studies submitted to the FDA were never
published, and among those that were, 11\% of primary endpoint
results were substantially different from the FDA's records.

The modification has two consequences. First, silence is
informationally ambiguous: the FDA cannot distinguish between a
firm that ran an informative experiment and got a bad result, and
a firm that ran a less informative experiment and got a bad
result. This prevents the complete unraveling that would occur
if the FDA could observe the experiment under silence. Second,
the FDA's approval probability after silence, $A(\emptyset)$,
becomes a single instrument that affects all firm types
simultaneously---unlike Min's type-specific approval
probabilities. This pooling constraint generates a new
tradeoff: $A(\emptyset)$ changes which menu option each firm
type prefers (the screening channel), potentially altering the
equilibrium assignment of experiments to types, but also lets
more bad drugs through when firms hide unfavorable evidence.


% ============================================================
\section{Model Setup}
% ============================================================

\subsection{Primitives}

There are two players: a pharmaceutical firm (the agent,
``he'') and the FDA (the principal, ``she'').

\paragraph{States and prior.}
The true quality of the drug is described by a state
$\omega \in \{G, B\}$, where $G$ denotes a good drug and $B$
denotes a bad drug. The common prior is
\begin{equation}
  \Pr(G) = p, \qquad \Pr(B) = 1 - p,
\end{equation}
where $p \in (0, 1/2)$. The restriction $p < 1/2$ means that,
absent any evidence, the drug is more likely bad than good.

\paragraph{Actions.}
The FDA chooses to approve ($A$) or reject ($D$).

\paragraph{Payoffs.}
The FDA receives 1 for a correct decision and 0 otherwise:
\begin{equation}\label{eq:fda_payoff}
\begin{aligned}
  U_P(A, G) &= 1, & U_P(A, B) &= 0, \\
  U_P(D, G) &= 0, & U_P(D, B) &= 1.
\end{aligned}
\end{equation}
The firm cares only about approval:
\begin{equation}\label{eq:firm_payoff}
  U_A(A, \cdot) = 1, \qquad U_A(D, \cdot) = 0.
\end{equation}
In matrix form:
\begin{center}
\begin{tabular}{c|cc}
         & $G$       & $B$       \\ \hline
  $A$    & $(1,\;1)$ & $(1,\;0)$ \\
  $D$    & $(0,\;0)$ & $(0,\;1)$
\end{tabular}
\end{center}
All primitives above are identical to Min (2023).


% ------------------------------------------------------------
\subsection{Types and Feasible Sets}
% ------------------------------------------------------------

The firm has a privately known type
$\theta_i \in \Theta = \{\theta_1, \theta_2\}$, with
$\Pr(\theta_1) = t$. Type determines the set of feasible
experiments $S_i$. Following Min (2023):

\begin{assumption}[Feasible sets]\label{as:feasible}
$S_2 \subset S_1$, and both $S_1$ and $S_2$ satisfy:
\begin{enumerate}
  \item[(i)] Each $S_i$ is a closed, convex subset of
    $\{(x, y) \in [0,1]^2 : y \geq 1 - x\}$.
  \item[(ii)] For each $S_i$, the upper boundary
    $\bar{S}_i(x) := \max\{y : (x,y) \in S_i\}$ is
    continuous, concave, and strictly decreasing on its
    domain.
\end{enumerate}
\end{assumption}

Here $\theta_1$ is the ``big type'' (e.g., a large firm with
a complex, multi-endpoint trial offering more analytical
flexibility) and $\theta_2$ is the ``small type'' (e.g., a
small firm or rare-disease trial with fewer endpoints and less
flexibility). The inclusion $S_2 \subset S_1$ means that any
experiment available to the small type is also available to the
big type, but not vice versa.


% ------------------------------------------------------------
\subsection{Experiments}
% ------------------------------------------------------------

An experiment $\pi = (x, y) \in S_i$ produces a binary outcome
$j \in \{g, b\}$:
\begin{equation}\label{eq:experiment}
\begin{aligned}
  \Pr(g \mid G, \pi) &= y,
    &\qquad \Pr(b \mid G, \pi) &= 1 - y, \\
  \Pr(g \mid B, \pi) &= 1 - x,
    &\qquad \Pr(b \mid B, \pi) &= x.
\end{aligned}
\end{equation}
Here $y$ is the true positive rate and $x$ is the true
negative rate. The restriction $y \geq 1 - x$ ensures that $g$
is ``good news'' and $b$ is ``bad news.''

An experiment is \emph{convincing} if a good result is
sufficient for the FDA to approve under ex post optimal
behavior, i.e.,
\begin{equation}\label{eq:convincing}
  py \geq (1 - p)(1 - x).
\end{equation}
All definitions above are identical to Min (2023).


% ============================================================
\subsection{The Modification: Voluntary Disclosure with
  Unobservable Experiments}
\label{sec:modification}
% ============================================================

In Min (2023), the FDA observes both the experiment $\pi$ and
the outcome $j$. I modify this as follows.

After observing $j \in \{g, b\}$, the firm chooses a message
$m$ to send to the FDA:
\begin{itemize}
  \item \textbf{Disclose:} $m = (\pi, j)$. The firm reveals
    both the experiment conducted and the realized outcome.
    This is verifiable: if the firm claims $(\pi, g)$, the true
    outcome must be $g$.
  \item \textbf{Silence:} $m = \emptyset$. The firm reveals
    nothing.
\end{itemize}

\medskip
\noindent\textbf{Key assumption:} When the firm discloses, the
FDA observes $(\pi, j)$---exactly as in Min. When the firm
remains silent, the FDA observes \emph{only} $\emptyset$: it
does not learn which experiment was conducted, nor what the
outcome was.

\medskip
This asymmetry is the single modification to Min's model. It
captures the empirical reality that published trial results are
verifiable and detailed (the medical community sees the trial
design, endpoints, sample size, and results), while buried
trials are often invisible to regulators: many are never
registered, and even registered trials frequently lack posted
results.

All other primitives---states, prior, payoffs, type structure,
feasible sets, and the definition of experiments---are
unchanged from Min (2023).


% ============================================================
\subsection{Timeline}
\label{sec:timeline}
% ============================================================

\begin{enumerate}
  \item The FDA commits to a decision rule: approval
    probabilities $A(\pi, j)$ for each possible disclosed
    evidence $(\pi, j)$, and a single approval probability
    $A(\emptyset)$ for silence.
  \item The firm observes the FDA's rule and privately learns
    its type $\theta_i$ (i.e., learns $S_i$).
  \item The firm selects and conducts an experiment
    $\pi \in S_i$. The outcome $j \in \{g, b\}$ is realized
    and privately observed by the firm.
  \item The firm chooses $m \in \{(\pi, j), \emptyset\}$.
  \item The FDA observes $m$ and makes its decision according
    to the committed rule.
\end{enumerate}

\begin{remark}[Comparison with Min's timeline]
Steps 1--3 are identical to Min (2023). The difference is in
step~4: Min's firm has no disclosure choice (the FDA
automatically sees $(\pi, j)$), whereas here the firm chooses
whether to disclose or remain silent. Crucially, if the firm
stays silent, the FDA does not observe $\pi$---this is what
distinguishes the model from a simple addition of
$\emptyset$ to Min's observation space.
\end{remark}


% ============================================================
\subsection{Strategies}
\label{sec:strategies}
% ============================================================

\paragraph{Firm's strategy.}
The firm's strategy has two components:
\begin{enumerate}
  \item \textbf{Experiment choice:} a mapping from type and
    decision rule to an experiment,
    $\sigma: \Theta \to S_i$, where $\sigma(\theta_i) \in S_i$.
  \item \textbf{Disclosure decision:} for each experiment $\pi$
    and outcome $j$, a probability of disclosing,
    \begin{equation}\label{eq:delta}
    \begin{aligned}
      \delta_g(\pi) &:= \Pr(m = (\pi, g) \mid j = g, \pi)
        \in [0, 1], \\
      \delta_b(\pi) &:= \Pr(m = (\pi, b) \mid j = b, \pi)
        \in [0, 1].
    \end{aligned}
    \end{equation}
    The probability of silence is $1 - \delta_j(\pi)$.
\end{enumerate}

\paragraph{FDA's decision rule.}
The FDA commits to:
\begin{equation}\label{eq:fda_rule}
  A(\pi, g), \quad A(\pi, b) \in [0, 1]
  \quad \text{for each } \pi \in S_1,
  \qquad\text{and}\qquad
  A(\emptyset) \in [0, 1].
\end{equation}
Note that $A(\pi, g)$ and $A(\pi, b)$ can be
experiment-specific (the FDA sees $\pi$ when the firm
discloses), but $A(\emptyset)$ is a \textbf{single number}
applying to all silent firms regardless of type or experiment.
This is because, under silence, the FDA does not observe $\pi$.


% ============================================================
\section{Analysis of the Disclosure Subgame}
\label{sec:disclosure}
% ============================================================

I first establish a basic fact, then analyze the firm's
optimal disclosure strategy and derive posterior beliefs.

\begin{lemma}[Bad news]\label{lem:bad_news}
For any experiment $\pi = (x,y)$ with $y \geq 1 - x$ and
prior $p < 1/2$:
\[
  \mu(G \mid \pi, b)
  = \frac{p(1-y)}{p(1-y) + (1-p)x}
  < \frac{1}{2}.
\]
\end{lemma}

\begin{proof}
$\mu(G \mid \pi, b) < 1/2$ iff $p(1-y) < (1-p)x$.
Since $p < 1/2$, we have $p/(1-p) < 1$.
Since $y \geq 1-x$, we have $1-y \leq x$.
Therefore $p(1-y) \leq p \cdot x < (1-p) \cdot x$.
\end{proof}

Fix an experiment $\pi = (x, y)$ and the FDA's rule
$(A(\pi, g), A(\pi, b), A(\emptyset))$. I analyze the
firm's optimal disclosure strategy, then derive the FDA's
posterior beliefs.

\subsection{Firm's Optimal Disclosure}

The firm's payoff equals the approval probability. For each
outcome $j$, the firm compares the payoff from disclosing
versus silence.

\paragraph{Upon $j = g$:}
Disclosing yields $A(\pi, g)$; silence yields $A(\emptyset)$.
The firm discloses if and only if
$A(\pi, g) \geq A(\emptyset)$.

\paragraph{Upon $j = b$:}
Disclosing yields $A(\pi, b)$; silence yields $A(\emptyset)$.
The firm discloses if and only if
$A(\pi, b) \geq A(\emptyset)$.

\begin{lemma}[Equilibrium disclosure]\label{lem:disclosure}
For a firm using experiment $\pi$ under FDA rule
$(A(\pi, g), A(\pi, b), A(\emptyset))$:
\begin{enumerate}
  \item[(i)] If $A(\pi, j) > A(\emptyset)$, outcome $j$ is
    disclosed: $\delta_j^* = 1$.
  \item[(ii)] If $A(\pi, j) < A(\emptyset)$, outcome $j$ is
    hidden: $\delta_j^* = 0$.
  \item[(iii)] If $A(\pi, j) = A(\emptyset)$, the firm is
    indifferent: any $\delta_j^* \in [0, 1]$ is consistent.
\end{enumerate}
\end{lemma}

\begin{proof}
The firm's payoff from disclosing $j$ is $A(\pi, j)$; from
silence it is $A(\emptyset)$. The firm chooses whichever is
higher.
\end{proof}

\begin{remark}[Typical case]
In the mechanism design analysis below, we will show that in
any optimal mechanism $A(\pi_i, g) \geq A(\emptyset)$ and
$A(\pi_i, b) \leq A(\emptyset)$
(Lemmas~\ref{lem:wlog_Ag} and~\ref{lem:hiding_dominates}).
Under these conditions, the firm discloses good results and
hides bad results---the ``cherry-picking'' behavior that is
this paper's focus.
\end{remark}


% ------------------------------------------------------------
\subsection{Posterior Beliefs}
\label{sec:posteriors}
% ------------------------------------------------------------

\paragraph{Upon disclosed evidence $(\pi, j)$.}
Since disclosure is verifiable, the posteriors are the same as
in Min (2023):
\begin{align}
  \mu(G \mid \pi, g)
    &= \frac{py}{py + (1-p)(1-x)},
    \label{eq:post_g} \\
  \mu(G \mid \pi, b)
    &= \frac{p(1-y)}{p(1-y) + (1-p)x}.
    \label{eq:post_b}
\end{align}
These do not depend on the disclosure probabilities (the
$\delta$'s cancel by Bayes' rule).

\paragraph{Upon silence $m = \emptyset$.}
This is the key object. Silence pools over all types and
experiments that produce undisclosed outcomes. Suppose in
equilibrium, type $\theta_i$ chooses experiment
$\pi_i = (x_i, y_i)$ and follows the disclosure strategy in
Lemma~\ref{lem:disclosure} (disclose $g$, hide $b$).

The probability of silence under each state is:
\begin{align}
  \Pr(\emptyset \mid G)
    &= \sum_{i=1}^{2} \Pr(\theta_i) \cdot (1 - y_i),
    \label{eq:silence_G} \\
  \Pr(\emptyset \mid B)
    &= \sum_{i=1}^{2} \Pr(\theta_i) \cdot x_i.
    \label{eq:silence_B}
\end{align}
Here, $(1 - y_i)$ is the probability of a bad result when the
drug is good and the firm uses experiment $\pi_i$, and $x_i$ is
the probability of a bad result when the drug is bad.

Applying Bayes' rule:
\begin{equation}\label{eq:post_silence}
\boxed{
  \mu(G \mid \emptyset)
  = \frac{p \sum_{i} t_i (1 - y_i)}
         {p \sum_{i} t_i (1 - y_i)
          + (1-p) \sum_{i} t_i \, x_i},
}
\end{equation}
where $t_1 = t$ and $t_2 = 1 - t$.

\begin{remark}[Contrast with Feb~27 model]
In the Feb~27 model, the FDA observed $(\pi, \emptyset)$---it
knew which experiment was conducted. The silence posterior
reduced to $\mu(G \mid \pi, b)$, leading to complete
unraveling: silence was perfectly informative, $A(\emptyset) = 0$,
and cherry-picking was irrelevant.

Here, the FDA sees only $\emptyset$. The posterior
\eqref{eq:post_silence} depends on the \emph{equilibrium mix}
of experiments chosen by both types. It is a weighted average
across types and experiments, not a single experiment's bad
posterior. This is what prevents unraveling.
\end{remark}


\begin{remark}[Why unraveling fails]\label{rem:no_unravel}
In the classic Milgrom--Grossman unraveling argument, the
receiver can infer the sender's information from silence
because the receiver knows the sender's ``type'' (i.e., what
information the sender could have). Here, the FDA does not know
the firm's type $\theta_i$, so it cannot determine which
experiment produced the hidden bad result. Different experiments
have different $(x, y)$, so the informational content of ``bad
result from $\pi_1$'' differs from ``bad result from $\pi_2$.''
Pooling across types and experiments makes silence genuinely
ambiguous.
\end{remark}


% ------------------------------------------------------------
\subsection{FDA's Optimal Response}
\label{sec:fda_response}
% ------------------------------------------------------------

The FDA approves if and only if its posterior exceeds $1/2$
(since its payoff is 1 for correct decisions, 0 otherwise):

\paragraph{Response to disclosed $(\pi, g)$:}
The FDA approves if $\mu(G \mid \pi, g) \geq 1/2$, i.e., if
the experiment is convincing:\ $py \geq (1-p)(1-x)$.
Same as Min.

\paragraph{Response to disclosed $(\pi, b)$:}
Under $p < 1/2$ and $y \geq 1 - x$, we have
$\mu(G \mid \pi, b) < 1/2$ (Lemma~\ref{lem:bad_news}). So the
FDA always rejects: $A(\pi, b) = 0$. Same as Min.

\paragraph{Response to silence $\emptyset$:}
The FDA approves if
$\mu(G \mid \emptyset) \geq 1/2$.

\begin{lemma}[Silence posterior is below bad-news posterior
  of the best experiment]\label{lem:silence_bound}
For any equilibrium experiment choices $\pi_1, \pi_2$:
\[
  \mu(G \mid \emptyset)
  \leq \max_{i \in \{1,2\}} \mu(G \mid \pi_i, b)
  < \frac{1}{2}.
\]
Therefore, under ex post optimal FDA behavior:
$A(\emptyset) = 0$.
\end{lemma}

\begin{proof}
Since silence occurs only when $j = b$ (because $\delta_g = 1$),
the silence posterior is a weighted average of the type-specific
bad-result posteriors:
\[
  \mu(G \mid \emptyset)
  = \frac{\sum_i t_i \, p(1-y_i)}
         {\sum_i t_i [p(1-y_i) + (1-p)x_i]}
  = \frac{\sum_i t_i \, p(1-y_i)}
         {\sum_i t_i \, \Pr(b \mid \pi_i)}.
\]
Each term $\frac{p(1-y_i)}{p(1-y_i) + (1-p)x_i}
= \mu(G \mid \pi_i, b) < 1/2$ by Lemma~\ref{lem:bad_news}.
A weighted average of quantities each strictly below $1/2$
is strictly below $1/2$.
\end{proof}

\begin{remark}[The ex post optimality trap]
Under ex post optimal behavior, $A(\emptyset) = 0$, and the
model collapses back to Min (2023): the firm's payoff from any
experiment becomes $\Pr(g \mid \pi) \cdot A(\pi, g)$,
identical to Min. Cherry-picking is again irrelevant.

This is exactly the advisor's critique of the Feb~27 model.
The result is trivial because the FDA's response to silence is
``too easy.''

The resolution, following Min's approach, is to recognize that
the FDA \textbf{commits} to a decision rule \emph{before}
observing the firm's choice. As Min shows, the FDA sometimes
benefits from committing to ex post suboptimal actions (e.g.,
$A(\pi, b) > 0$) to screen firm types. Similarly, here the FDA
may benefit from committing to $A(\emptyset) > 0$---not
because silence is good news, but because this commitment
changes the firm's experiment-choice incentives.
\end{remark}


% ============================================================
\section{The Mechanism Design Problem}
\label{sec:mechanism}
% ============================================================

Following Min (2023), I formulate the FDA's problem as a
mechanism design problem using the revelation principle.

\subsection{General Menu}

The FDA offers a menu of options, one intended for each type,
together with a silence approval probability:
\begin{equation}\label{eq:menu_full}
  \mathcal{M} = \left\{
    \left(\pi_1, A_g^1, A_b^1\right),\;
    \left(\pi_2, A_g^2, A_b^2\right);\;
    A(\emptyset)
  \right\},
\end{equation}
where:
\begin{itemize}
  \item $\pi_i = (x_i, y_i) \in S_i$ is the experiment
    intended for type $\theta_i$.
  \item $A_g^i \in [0, 1]$ is the approval probability when
    the firm discloses a good result from experiment $\pi_i$.
  \item $A_b^i \in [0, 1]$ is the approval probability when
    the firm discloses a bad result from experiment $\pi_i$.
  \item $A(\emptyset) \in [0, 1]$ is the approval probability
    upon silence---\textbf{shared across all types}.
\end{itemize}

By Lemma~\ref{lem:disclosure}, type $\theta_i$'s disclosure
behavior depends on comparing $A_g^i$, $A_b^i$ with
$A(\emptyset)$:
\begin{itemize}
  \item If $A_b^i > A(\emptyset)$: firm discloses $b$
    (the FDA sees $(\pi_i, b)$ and approves with prob
    $A_b^i$). This is a \textbf{disclosure regime} for type
    $i$, identical to Min.
  \item If $A_b^i \leq A(\emptyset)$: firm hides $b$
    (silence; FDA approves with prob $A(\emptyset)$). This
    is the \textbf{hiding regime}---the paper's focus.
\end{itemize}

The next two lemmas show that the menu
\eqref{eq:menu_full} can be simplified.

\begin{lemma}[WLOG: $A_g^i \geq A(\emptyset)$]
\label{lem:wlog_Ag}
In any optimal mechanism, $A_g^i \geq A(\emptyset)$ for each
type~$i$.
\end{lemma}

\begin{proof}
Suppose $A_g^i < A(\emptyset)$. Then by
Lemma~\ref{lem:disclosure}, type~$i$ hides both $g$ and $b$
(always silent), earning $A(\emptyset)$ regardless of
experiment or outcome. Replace $A_g^i$ with $A(\emptyset)$.

\emph{Firm payoffs are unchanged:} type $i$'s payoff from
option~$i$ remains $A(\emptyset)$ (the firm is now indifferent
about disclosing $g$; WLOG assume $\delta_g = 1$, but payoff
is the same either way since $A_g^i = A(\emptyset)$). Any
other type deviating to option~$i$ also gets $A(\emptyset)$
from option~$i$, same as before.

\emph{IC constraints are preserved:} all firm payoffs from
all options are unchanged.

\emph{FDA payoff is unchanged:} the FDA approves with
probability $A(\emptyset)$ in all events involving type~$i$,
same as before.

So the modified mechanism is equivalent. From
$A_g^i = A(\emptyset)$, the FDA can weakly improve by
increasing $A_g^i$ above $A(\emptyset)$ whenever the assigned
experiment is convincing ($\Delta_g(\pi_i) > 0$), since this
strictly improves the FDA's payoff from type~$i$'s
good-result events.
\end{proof}

\begin{lemma}[Hiding dominates disclosure of bad results]
\label{lem:hiding_dominates}
Fix an optimal mechanism with $A_g^i < 1$ for type~$i$.
If $A_b^i > A(\emptyset)$ (so type~$i$ discloses $b$),
the FDA can weakly improve by switching to hiding:
set $A_b^i' = 0$ and raise $A_g^i$ by
$\Delta := (A_b^i - A(\emptyset)) \cdot
\Pr(b \mid \pi_i) / \Pr(g \mid \pi_i)$.
\end{lemma}

\begin{proof}
\emph{Step 1: Firm payoffs are unchanged.}
Under the original mechanism, type~$i$'s payoff from option~$i$
is $A_g^i \Pr(g|\pi_i) + A_b^i \Pr(b|\pi_i)$ (disclosing both
$g$ and $b$). Under the modified mechanism, it is
$(A_g^i + \Delta) \Pr(g|\pi_i) + A(\emptyset) \Pr(b|\pi_i)$.
By construction of $\Delta$, these are equal. Any other type
deviating to option~$i$ faces the same payoff calculation, so
all firm payoffs from option~$i$ are unchanged. IC is
preserved.

\emph{Step 2: FDA payoff improves.}
The FDA's payoff from type~$i$ is
$EU^P_i = (1-p) + \Delta_g(\pi_i) A_g^i
+ \Delta_b(\pi_i) A_{\text{eff}}^i$, where
$\Delta_g := py_i - (1-p)(1-x_i)$,
$\Delta_b := p(1-y_i) - (1-p)x_i$, and $A_{\text{eff}}^i$ is
the effective bad-result approval ($A_b^i$ under disclosure,
$A(\emptyset)$ under hiding). The change is:
\[
  \Delta_g(\pi_i) \cdot \Delta
  + \Delta_b(\pi_i) \cdot (A(\emptyset) - A_b^i).
\]
Substituting $\Delta$ and simplifying, the net change equals:
\[
  \frac{(A_b^i - A(\emptyset))}{\Pr(g|\pi_i)}
  \left[\Delta_g(\pi_i) \Pr(b|\pi_i)
    - \Delta_b(\pi_i) \Pr(g|\pi_i)\right].
\]
The bracketed expression expands to
$2p(1-p)(x_i + y_i - 1) \geq 0$ (since $y_i \geq 1 - x_i$
by the information constraint), and
$A_b^i - A(\emptyset) > 0$ by assumption. So the FDA's payoff
weakly improves (strictly when $x_i + y_i > 1$).
\end{proof}

\begin{remark}[Scope of Lemma~\ref{lem:hiding_dominates}]
The lemma requires $A_g^i < 1$ (so there is room to raise
$A_g^i$ to compensate IC). When $A_g^i = 1$, the compensation
argument does not directly apply. However, under $A_g^i = 1$
the firm already receives near-maximal approval; further
raising the bad-result approval $A_b^i$ above $A(\emptyset)$
would require the FDA to approve type~$i$ even more
generously, which is costly. A full characterization of this
corner case requires solving the optimization problem for
specific feasible sets.
\end{remark}

Combining Lemmas~\ref{lem:wlog_Ag}
and~\ref{lem:hiding_dominates}, I henceforth focus on
mechanisms satisfying $A_g^i \geq A(\emptyset) \geq A_b^i$.
Under these conditions, the firm discloses $g$ and hides $b$,
and $A_b^i$ is payoff-irrelevant (the effective bad-result
approval is $A(\emptyset)$). The menu simplifies to:
\begin{equation}\label{eq:menu}
  \mathcal{M} = \left\{
    \left(\pi_1, A_g^1\right),\;
    \left(\pi_2, A_g^2\right);\;
    A(\emptyset)
  \right\}.
\end{equation}

\begin{remark}[Comparison with Min's menu]
\label{rem:menu_comparison}
In Min (2023), the FDA always observes $(\pi, j)$, so the menu
is $\{(\pi_1, A_g^1, A_b^1), (\pi_2, A_g^2, A_b^2)\}$---6
instruments (plus experiments). Each type has its own
bad-result approval probability $A_b^i$.

In my model, Lemma~\ref{lem:hiding_dominates} shows that the
FDA (generically) prefers the firm to hide bad results. The
effective instruments are $(\pi_1, \pi_2, A_g^1, A_g^2,
A(\emptyset))$---5 instruments. The type-specific $A_b^i$
collapses to a single pooled $A(\emptyset)$: the FDA has one
fewer screening tool.
\end{remark}


% ------------------------------------------------------------
\subsection{Firm's Expected Payoff}
\label{sec:firm_payoff}
% ------------------------------------------------------------

Under the simplified menu \eqref{eq:menu} with
$A_g^i \geq A(\emptyset)$, the firm discloses $g$ and hides
$b$ (Lemma~\ref{lem:disclosure}). The firm's expected payoff
from choosing option $(\pi_i, A_g^i)$ is:
\begin{equation}\label{eq:EU_A}
  EU^A(\pi_i, A_g^i, A(\emptyset))
  = A_g^i \cdot \Pr(g \mid \pi_i)
    + A(\emptyset) \cdot \Pr(b \mid \pi_i),
\end{equation}
where
\begin{align}
  \Pr(g \mid \pi_i) &= p \, y_i + (1-p)(1-x_i),
    \label{eq:prob_g} \\
  \Pr(b \mid \pi_i) &= p(1-y_i) + (1-p) \, x_i.
    \label{eq:prob_b}
\end{align}

\begin{remark}[Comparison with Min's firm payoff]
In Min (2023):
$EU^A_{\text{Min}} = A_g^i \cdot \Pr(g \mid \pi_i)
  + A_b^i \cdot \Pr(b \mid \pi_i)$.

My model replaces the type-specific $A_b^i$ with the
pooled $A(\emptyset)$. When $A(\emptyset) = 0$, the two
payoffs coincide (both reduce to
$A_g^i \cdot \Pr(g \mid \pi_i)$).

When $A(\emptyset) > 0$, the firm gets a ``floor'' payoff
$A(\emptyset) \cdot \Pr(b \mid \pi_i)$ from hiding bad
results. This floor is the same regardless of type---unlike
Min, where the FDA can set different floors for different
types.
\end{remark}


% ------------------------------------------------------------
\subsection{Incentive Compatibility}
\label{sec:ic}
% ------------------------------------------------------------

The IC constraints ensure that each type prefers its intended
option. Type $\theta_1$ (big) must prefer option 1 over
option 2:
\begin{equation}\label{eq:IC1}
  A_g^1 \cdot \Pr(g \mid \pi_1)
    + A(\emptyset) \cdot \Pr(b \mid \pi_1)
  \geq
  A_g^2 \cdot \Pr(g \mid \pi_2)
    + A(\emptyset) \cdot \Pr(b \mid \pi_2).
\end{equation}
Type $\theta_2$ (small) must prefer option 2 over option 1.
If $\pi_1 \in S_1 \setminus S_2$, this is automatically
satisfied since $\theta_2$ cannot conduct $\pi_1$. If
$\pi_1 \in S_2$, then:
\begin{equation}\label{eq:IC2}
  A_g^2 \cdot \Pr(g \mid \pi_2)
    + A(\emptyset) \cdot \Pr(b \mid \pi_2)
  \geq
  A_g^1 \cdot \Pr(g \mid \pi_1)
    + A(\emptyset) \cdot \Pr(b \mid \pi_1).
\end{equation}

\begin{remark}[How $A(\emptyset)$ changes IC]
Rewrite IC$_1$ \eqref{eq:IC1} as:
\[
  A_g^1 \Pr(g \mid \pi_1) - A_g^2 \Pr(g \mid \pi_2)
  \geq A(\emptyset)
    \left[\Pr(b \mid \pi_2) - \Pr(b \mid \pi_1)\right].
\]
Since $\Pr(b \mid \pi) = 1 - \Pr(g \mid \pi)$, the
right-hand side equals
$A(\emptyset)[\Pr(g \mid \pi_1) - \Pr(g \mid \pi_2)]$.

If experiment $\pi_1$ is more likely to produce $g$ than
$\pi_2$ (i.e., $\Pr(g \mid \pi_1) > \Pr(g \mid \pi_2)$),
then raising $A(\emptyset)$ \textbf{tightens} IC$_1$: the big
type's ``insurance'' from silence is greater when using the
experiment more likely to produce $b$ (which is $\pi_2$),
making deviation to option~2 more attractive.

Conversely, if $\pi_1$ produces $g$ less often than $\pi_2$,
raising $A(\emptyset)$ \textbf{relaxes} IC$_1$. This
interaction between $A(\emptyset)$ and experiment choice is
absent from Min.
\end{remark}


% ------------------------------------------------------------
\subsection{FDA's Posterior under Silence (Consistency)}
\label{sec:consistency}
% ------------------------------------------------------------

The FDA's committed $A(\emptyset)$ must be consistent with
Bayesian beliefs in equilibrium. If the FDA commits to
$A(\emptyset) > 0$, then in equilibrium both types hide bad
results. The silence posterior is
\eqref{eq:post_silence}:
\begin{equation}\label{eq:consistency}
  \mu(G \mid \emptyset)
  = \frac{p \left[t(1-y_1) + (1-t)(1-y_2)\right]}
         {p \left[t(1-y_1) + (1-t)(1-y_2)\right]
          + (1-p)\left[t \, x_1 + (1-t) x_2\right]}.
\end{equation}

By Lemma~\ref{lem:silence_bound}, this posterior is always
below $1/2$. So from an ex post perspective, the FDA should
reject. But the FDA is committing \emph{ex ante}---before
seeing the firm's choice. This is the same logic as Min's
commitment to $A_b^i > 0$: the FDA sacrifices ex post
optimality to improve ex ante screening.

The FDA's expected cost of $A(\emptyset) > 0$ is:
\begin{equation}\label{eq:fda_cost_silence}
  \text{Cost} = A(\emptyset) \cdot \Pr(\emptyset)
    \cdot \left[1 - \mu(G \mid \emptyset)\right],
\end{equation}
which is the probability of approving a bad drug through the
silence channel.


% ============================================================
\subsection{FDA's Optimization Problem}
\label{sec:fda_problem}
% ============================================================

The FDA maximizes its expected payoff over the menu
$\mathcal{M} = \{(\pi_1, A_g^1), (\pi_2, A_g^2);
A(\emptyset)\}$, subject to IC constraints and feasibility.

The FDA's expected payoff from type $\theta_i$ choosing option
$(\pi_i, A_g^i)$ and hiding bad results (disclose $g$, hide
$b$). The FDA gets payoff 1 for a correct decision and 0
otherwise:

\begin{itemize}
  \item \textbf{State $G$, outcome $g$} (prob $p \, y_i$):
    firm discloses, FDA approves with prob $A_g^i$.
    FDA payoff: $A_g^i \cdot 1 + (1-A_g^i) \cdot 0 = A_g^i$.
  \item \textbf{State $G$, outcome $b$} (prob $p(1-y_i)$):
    firm silent, FDA approves with prob $A(\emptyset)$.
    FDA payoff: $A(\emptyset) \cdot 1 + (1-A(\emptyset))
    \cdot 0 = A(\emptyset)$.
    (Correct to approve since state is $G$.)
  \item \textbf{State $B$, outcome $g$} (prob $(1-p)(1-x_i)$):
    firm discloses, FDA approves with prob $A_g^i$.
    FDA payoff: $A_g^i \cdot 0 + (1-A_g^i) \cdot 1
    = 1 - A_g^i$.
    (Incorrect to approve since state is $B$.)
  \item \textbf{State $B$, outcome $b$} (prob $(1-p) x_i$):
    firm silent, FDA approves with prob $A(\emptyset)$.
    FDA payoff: $A(\emptyset) \cdot 0 + (1-A(\emptyset))
    \cdot 1 = 1 - A(\emptyset)$.
    (Correct to reject since state is $B$.)
\end{itemize}

Therefore:
\begin{equation}\label{eq:EU_P}
\begin{aligned}
  EU^P_i
  &= p \, y_i \, A_g^i
     + p(1-y_i) \, A(\emptyset)
     + (1-p)(1-x_i)(1 - A_g^i)
     + (1-p) \, x_i \, (1 - A(\emptyset)).
\end{aligned}
\end{equation}

Simplifying:
\begin{equation}\label{eq:EU_P_simple}
\begin{aligned}
  EU^P_i
  &= p \, y_i \, A_g^i
     + p(1-y_i) \, A(\emptyset)
     + (1-p)(1-x_i) - (1-p)(1-x_i) A_g^i
     + (1-p) x_i - (1-p) x_i \, A(\emptyset) \\
  &= (1-p) + \left[p \, y_i - (1-p)(1-x_i)\right] A_g^i
     + \left[p(1-y_i) - (1-p) x_i\right] A(\emptyset).
\end{aligned}
\end{equation}

Define, for each experiment $\pi_i = (x_i, y_i)$:
\begin{align}
  \Delta_g(\pi_i)
    &:= p \, y_i - (1-p)(1-x_i),
    \label{eq:delta_g_def} \\
  \Delta_b(\pi_i)
    &:= p(1-y_i) - (1-p) x_i.
    \label{eq:delta_b_def}
\end{align}
Then:
\begin{equation}\label{eq:EU_P_compact}
  EU^P_i = (1-p) + \Delta_g(\pi_i) \, A_g^i
           + \Delta_b(\pi_i) \, A(\emptyset).
\end{equation}

\begin{remark}[Interpretation of $\Delta_g$ and $\Delta_b$]
$\Delta_g(\pi_i) \geq 0$ iff the experiment is convincing
\eqref{eq:convincing}: a good result is more likely under $G$
than under $B$ (weighted by priors). For convincing experiments,
$A_g^i > 0$ helps the FDA.

$\Delta_b(\pi_i) < 0$ always (by Lemma~\ref{lem:bad_news}):
a bad result is more likely under $B$ than under $G$.
So $A(\emptyset) > 0$ \textbf{hurts} the FDA---it means
approving when the evidence (silence = hidden bad result)
points toward $B$.

This is exactly like Min's tradeoff with $A_b^i$: the FDA
sometimes commits to costly actions ($A_b > 0$ in Min,
$A(\emptyset) > 0$ here) to improve screening.
\end{remark}

The FDA's overall problem is:
\begin{equation}\label{eq:fda_max}
  \max_{\substack{\pi_1, \pi_2, A_g^1, A_g^2, A(\emptyset)}}
  \; t \cdot EU^P_1 + (1-t) \cdot EU^P_2
\end{equation}
subject to:
\begin{enumerate}
  \item $\pi_i \in S_i$ for $i = 1, 2$,
  \item $A_g^1, A_g^2, A(\emptyset) \in [0, 1]$,
  \item IC constraints \eqref{eq:IC1} and \eqref{eq:IC2}.
\end{enumerate}

\begin{remark}[Comparison with Min's problem]
Min's FDA problem has the same structure but with separate
$A_b^1, A_b^2$ instead of a shared $A(\emptyset)$. Min has
6 choice variables ($\pi_1, \pi_2, A_g^1, A_g^2, A_b^1,
A_b^2$); I have 5 ($\pi_1, \pi_2, A_g^1, A_g^2, A(\emptyset)$).
The reduction from 6 to 5 is exactly the pooling constraint:
$A_b^1 = A_b^2 = A(\emptyset)$.

If the pooling constraint is not binding (i.e., the optimal
$A_b^1 = A_b^2$ in Min's problem), then my model's solution
coincides with Min's. If the constraint is binding, my model
produces genuinely different results.
\end{remark}


% ============================================================
\section{Benchmark: Ex Post Optimality}
\label{sec:benchmark}
% ============================================================

As a benchmark, consider the case where the FDA plays ex post
optimally---it sets $A(\emptyset) = 0$ because silence is bad
news (Lemma~\ref{lem:silence_bound}).

\begin{proposition}[Benchmark: trivial cherry-picking]
\label{prop:benchmark}
If the FDA responds ex post optimally to each observation,
then $A(\emptyset) = 0$, the firm's disclosure choice is
payoff-irrelevant, and the equilibrium is identical to Min
(2023).
\end{proposition}

\begin{proof}
With $A(\emptyset) = 0$, the firm's payoff is
$A_g^i \cdot \Pr(g \mid \pi_i)$ regardless of disclosure
behavior. This is exactly Min's firm payoff with $A_b^i = 0$.
\end{proof}

This confirms the advisor's observation: the Feb~27 result is
trivial because it only considers ex post optimal FDA behavior.
The interesting question is whether the FDA benefits from
\emph{committing} to $A(\emptyset) > 0$.


% ============================================================
\section{When Does Positive $A(\emptyset)$ Help the FDA?}
\label{sec:positive_A}
% ============================================================

\subsection{The Effect of $A(\emptyset)$ on Experiment Choice}

Raising $A(\emptyset)$ has two effects on the firm's
experiment choice:

\paragraph{Direct effect (bad for FDA):}
For a fixed experiment $\pi_i$, a higher $A(\emptyset)$
increases the firm's payoff through the $A(\emptyset) \cdot
\Pr(b \mid \pi_i)$ term. The firm gets approved more often
even with bad results. This is costly for the FDA because
$\Delta_b < 0$.

\paragraph{Screening effect (the real channel):}
$A(\emptyset)$ changes which \emph{menu option} each type
prefers. The firm's payoff from option $i$ is:
\[
  EU^A_i = A_g^i \cdot \Pr(g \mid \pi_i)
         + A(\emptyset) \cdot \Pr(b \mid \pi_i)
         = A(\emptyset)
           + (A_g^i - A(\emptyset)) \cdot \Pr(g \mid \pi_i).
\]
A higher $A(\emptyset)$ compresses the payoff difference
between options with high vs.\ low $\Pr(g \mid \pi_i)$. This
can change the big type's preference over which option to
select, altering the binding IC constraint and therefore the
FDA's optimal menu design.

\begin{lemma}[Experiment choice]\label{lem:exp_choice}
For a given menu option $(A_g^i, A(\emptyset))$ with
$A_g^i > A(\emptyset)$, the firm's optimal experiment within
$S_i$ solves:
\[
  \max_{\pi \in S_i} \; (A_g^i - A(\emptyset))
    \cdot \Pr(g \mid \pi) + A(\emptyset).
\]
Since $A(\emptyset)$ is a constant, this reduces to:
\[
  \max_{\pi \in S_i} \; \Pr(g \mid \pi)
  = \max_{(x,y) \in S_i} \; py + (1-p)(1-x).
\]
The optimal experiment does \textbf{not} depend on
$A(\emptyset)$ (given $A_g^i > A(\emptyset)$).
\end{lemma}

\begin{remark}[Within-option neutrality---same as Min]
The firm's optimal experiment within a given menu option does
not depend on $A(\emptyset)$: the firm always maximizes
$\Pr(g \mid \pi)$. This is \textbf{identical to Min (2023)},
where the firm's optimal experiment within a menu option does
not depend on $A_b^i$ either (since
$EU^A_{\text{Min}} = A_b^i + (A_g^i - A_b^i) \Pr(g|\pi)$,
and the firm maximizes $\Pr(g|\pi)$ regardless of $A_b^i$).

In both Min and this model, the interesting effects of the
``bad-outcome approval probability'' ($A_b^i$ in Min,
$A(\emptyset)$ here) operate through the
\textbf{screening channel}: they affect which menu option each
type prefers (IC constraints), and therefore which experiment
each type ends up conducting in equilibrium. The difference is
that Min can set $A_b^1 \neq A_b^2$ to screen types
independently, while here $A(\emptyset)$ is a single pooled
instrument affecting both types simultaneously.
\end{remark}


% ------------------------------------------------------------
\subsection{The Screening Tradeoff}
\label{sec:screening_tradeoff}
% ------------------------------------------------------------

The FDA's core tradeoff with $A(\emptyset)$ operates through
IC constraints:

\paragraph{Experiment assignment.}
Following Min, consider assigning $\pi_1 \in S_1 \setminus S_2$
(big type does an experiment small type cannot) and
$\pi_2 \in S_2$ (small type's experiment). Then IC$_2$ is
automatic (small type cannot deviate), and only IC$_1$ binds.

IC$_1$ becomes:
\begin{equation}\label{eq:IC1_expanded}
  A_g^1 \Pr(g \mid \pi_1) + A(\emptyset) \Pr(b \mid \pi_1)
  \geq
  A_g^2 \Pr(g \mid \pi_2) + A(\emptyset) \Pr(b \mid \pi_2).
\end{equation}

Rearranging:
\begin{equation}\label{eq:IC1_rearranged}
  A_g^1 \Pr(g \mid \pi_1) - A_g^2 \Pr(g \mid \pi_2)
  \geq A(\emptyset) \left[\Pr(b \mid \pi_2)
    - \Pr(b \mid \pi_1)\right].
\end{equation}

\paragraph{Case 1:}
$\Pr(g \mid \pi_1) > \Pr(g \mid \pi_2)$ (equivalently,
$\Pr(b \mid \pi_2) > \Pr(b \mid \pi_1)$). The big type's
experiment produces good results more often. Then the RHS
of \eqref{eq:IC1_rearranged} is positive when
$A(\emptyset) > 0$: raising $A(\emptyset)$
\textbf{tightens} IC$_1$.

Intuition: when $A(\emptyset) > 0$, the big type gets
insurance against bad results. Since the small type's
experiment produces bad results more often, deviating to the
small type's option gives the big type \emph{more} insurance
payoffs, making deviation more attractive. The FDA must then
distort $A_g^1$ or $A_g^2$ more to maintain IC, which is
costly.

\paragraph{Case 2:}
$\Pr(g \mid \pi_1) < \Pr(g \mid \pi_2)$. The big type's
assigned experiment is less likely to produce good results
(perhaps because it is more informative/stringent). Then
raising $A(\emptyset)$ \textbf{relaxes} IC$_1$.

Intuition: the big type's experiment produces bad results
more often, so the big type benefits more from the silence
insurance. This makes the big type less tempted to deviate to
the small type's option.

\medskip
In Min (2023), the FDA can handle these cases by independently
adjusting $A_b^1$ and $A_b^2$. Here, the FDA has only one
instrument $A(\emptyset)$ affecting both IC constraints
simultaneously. This is the source of genuinely new results.


% ============================================================
\section{Discussion}
\label{sec:discussion}
% ============================================================

\subsection{Connection to Herresthal (2022)}

Herresthal (2022) studies hidden testing and selective
disclosure in a sequential-testing model. Her key result is
that hidden testing can yield a strict Pareto improvement over
observable testing when the conflict of interest between the
agent and DM is large. The mechanism is that hidden testing
encourages the agent to acquire \emph{more} information (since
bad results can be concealed), and the extra information
sometimes reveals that the product is bad---benefiting both
parties.

My model has an analogous mechanism operating through
experiment choice rather than testing duration. When
$A(\emptyset) > 0$, the firm is less afraid of bad results
and may be willing to accept a more informative experiment
assignment from the FDA (one with higher $x$ and $y$). This
can benefit both the firm (higher expected approval) and the
FDA (more informative evidence when disclosed).

The key difference is that Herresthal compares two regimes
(observable vs.\ hidden testing), while I study optimal
mechanism design within the hidden-disclosure regime,
following Min's screening framework.

\subsection{Policy Implications}

\paragraph{Pre-registration as shrinking the gap between
  regimes.}
Mandatory trial pre-registration (e.g., ClinicalTrials.gov)
effectively moves the model toward Min: the FDA observes
$\pi$ even under silence, so
$A(\emptyset) \to A(\pi, \emptyset)$ becomes
experiment-specific. This breaks the pooling constraint and
allows more tailored screening. The model predicts that
pre-registration is most valuable when firm types are similar
(small $S_1 \setminus S_2$), and least valuable---or even
harmful---when types are very different.

\paragraph{Audit as an intermediate regime.}
If the FDA can audit silent firms with probability $\alpha$
(discovering $\pi$ and $j$), the model interpolates between
the full-pooling regime ($\alpha = 0$, my model) and full
transparency ($\alpha = 1$, Min). Optimal audit intensity
balances the benefit of screening precision against the cost
of discouraging informative experiment choices.

\paragraph{Empirical prediction.}
The model predicts that in therapeutic areas with high type
heterogeneity (e.g., large vs.\ small firms, complex vs.\
simple trials), stricter disclosure enforcement may lead to
\emph{less} informative experiments---firms choose ``safer''
endpoints with lower informativeness to avoid the risk of
visible failure. This is consistent with the finding that
transparency mandates reduce innovative trial investment
(MSOM/INFORMS, 2023).


% ============================================================
\section{Next Steps}
\label{sec:next_steps}
% ============================================================

\begin{enumerate}
  \item \textbf{Characterize optimal $A(\emptyset)$.}
    Solve the FDA's problem \eqref{eq:fda_max} for specific
    feasible-set structures (e.g., linear or circular $S_i$
    boundaries) and characterize when the optimal
    $A(\emptyset) > 0$.

  \item \textbf{Compare with Min's solution.}
    Show formally when the pooling constraint
    ($A_b^1 = A_b^2 = A(\emptyset)$) is binding, i.e., when
    Min's optimal $A_b^1 \neq A_b^2$ and the disclosure model
    produces different results.

  \item \textbf{Welfare comparison.}
    Compare the FDA's payoff under (a) full transparency
    (Min), (b) hidden disclosure with optimal $A(\emptyset)$
    (this model), and (c) hidden disclosure with ex post
    optimal FDA ($A(\emptyset) = 0$, the trivial case).

  \item \textbf{Incorporate Herresthal-type results.}
    Show conditions under which the firm's experiment choice
    under hidden disclosure with optimal commitment is more
    informative than under full transparency, yielding a
    potential Pareto improvement.
\end{enumerate}


\end{document}
